\section{共同算术目标、量词与路线分工}\label{sec:arithmetic}
\subsection{终点与辅助量}
记 $\pi_2(y)=\#\{p\le y:p,p+2\text{ 均为素数}\}$。终点是 $\pi_2(y)\to\infty$，不是有限区间内找到更多素数对。一个常用加权辅助量是
\[
S_2(y)=\sum_{n\le y}\Lambda(n)\Lambda(n+2),
\]
其中 $\Lambda$ 是 von Mangoldt 函数。若证明某个正主项规模的渐近式，例如 $S_2(y)\sim c y$ 且 $c>0$，则可推出孪生素数无穷；这里的渐近式只是充分条件，本文没有证明它。

这一蕴含需要排除真素数幂的贡献。至少一个坐标为 $p^k$、$k\ge2$ 的整数对有 $O(\sqrt y\log y)$ 个，每项权重至多 $O(\log^2 y)$，总贡献为 $o(y)$。若真素数对只有有限多个，剩余加权贡献也不可能形成正的线性主项。这个初等核查提醒我们：即使辅助量是标准形式，其与终点的连接也应写出来，而不能只借助命名。

\subsection{岛①：局部筛因子为何还不是素数对计数}
令 $y\ge2$，$P_y=\prod_{p\le y}p$、$s_y(n)=\1_{(n,P_y)=1}$、$\rho_y=\prod_{p\le y}(1-1/p)$。
对固定有限非空的互异整数移位集 $\mathcal H=\{h_1,\ldots,h_r\}\subset\ZZ$，记 $\nu_p(\mathcal H)$ 为其模 $p$ 的不同剩余类数。
CRT 给出完整周期的精确恒等式
\[
C_y(\mathcal H):=\frac1{P_y}\sum_{n\bmod P_y}\prod_{h\in\mathcal H}s_y(n+h)
=\prod_{p\le y}\left(1-\frac{\nu_p(\mathcal H)}p\right).
\]
证明只需逐素数数出 $p-\nu_p(\mathcal H)$ 个允许类，再由 CRT 相乘。
若 $\mathcal H$ 可容许，即每个 $\nu_p<p$，则对足够大的 $p$ 有 $\nu_p=r$，
归一化局部因子 $(1-r/p)(1-1/p)^{-r}=1+O_r(p^{-2})$，因此乘积收敛。
特别对 $\{0,2\}$，
\[
\frac{C_y(\{0,2\})}{\rho_y^2}\longrightarrow 2C_2,\qquad
C_2=\prod_{p>2}\frac{p(p-2)}{(p-1)^2}>0.
\]
这是经典局部计算，不是本研究新得到的素数对渐近式\cite{archive1}。

将 $\min\mathcal H$ 平移为零并记 $h_{\max}=\max\mathcal H$，令 $X\to\infty$。
若 $Y_X=\sqrt{X+h_{\max}}$ 且 $Y_X<n\le X$，则所有 $n+h$ 均大于 $Y_X$、不超过 $Y_X^2$。
它们通过 $Y_X$ 以下素数筛，当且仅当它们都是素数。这是\emph{逐点识别}，没有估计通过筛的数量。
另一方面，对完整的 $1\le n\le X$ 窗口，周期分割只给出
\[
\left|\sum_{1\le n\le X}\prod_{h\in\mathcal H}s_y(n+h)-XC_y(\mathcal H)\right|
\le P_y+1.
\]
因为完整周期的贡献精确，余下不足一个周期的部分误差至多 $P_y$，非整数端点再付 $1$。
当 $P_y=o(X\rho_y^r)$ 时，上述周期截断界保证归一化误差趋于零；
而在精确识别尺度 $y=Y_X$ 上，素数定理给出 $\log P_{Y_X}\sim\sqrt X$，
故这一误差上界不能提供所需的相对误差控制。
这只是两种初等论证适用尺度不相接，不是所有筛方法的不可能定理。

\subsection{七岛不是单一串行证明}
\begin{table}[ht]
\centering\small
\caption{七岛的数学分工。表中“尚需”不是对全部可能路线的必要性断言。}
\begin{tabularx}{\textwidth}{c l X}
\toprule
岛 & 对象或任务 & 在本框架中尚需辨认的转移\\
\midrule
1 & 算术基础 & 从筛法、相关量或编码出发，保留原始计数解释\\
2 & 解析消去 & 从局部结构到完整实际系数的有符号节省\\
3 & 固定原子 & 指定移位、指定相位与固定 packet 的逐点控制\\
4 & 实际出现项／H1 & 完整载体、来源对应、重数与逐系数守恒\\
5 & 非自治动力学 & 从测度命题到指定算术初始点的命中结论\\
6 & 孪生素数终点 & 无穷多个真实素数对，而非模型或有限证书\\
7 & 几何辅助提升 & 保时对应与事件读出，不额外获得算术典型性\\
\bottomrule
\end{tabularx}
\end{table}

\subsection{每次转移携带的契约}
为比较来自不同阶段的公式，本文把一个来源契约记为
\[
\mathcal O=(\text{系数},h_0,\text{支撑与顺序},\text{参数范围},
\text{量词},\text{归一化},\text{损失预算}).
\]
这不是一个新增数学对象，而是防止误用定理的记录方式。两个公式符号相似，并不表示它们的 $\mathcal O$ 相同。例如，任意大的 CRT 起点不等于给定物理区间内的起点；在一组完整壳层上的前缀不等于任意素数删选；重新计算局部行能量也不等于压缩原有归一化矩阵。

本文使用四种证据表述。第一种是可以直接展开验证的有限恒等式；第二种是带有完整假设与证明的统一命题；第三种是明确保留未证前提的条件推论；第四种是有限数值诊断。数值诊断不承担无限量词，条件推论不会因它的程序检查通过而变成无条件定理。本文后续的每个命题都应按这一层次阅读，而非按它在研究档案中出现的先后阅读。
